\documentclass[12pt,a4paper]{article}

\usepackage[a4paper,margin=1in]{geometry}
\usepackage[T1]{fontenc}
\usepackage[utf8]{inputenc}
\usepackage{xcolor}
\usepackage{hyperref}
\usepackage{tabularx}
\usepackage{booktabs}
\usepackage{array}
\usepackage{enumitem}
\usepackage{titlesec}
\usepackage{tcolorbox}
\usepackage{fancyhdr}
\usepackage{setspace}
\usepackage[table]{xcolor}

\hypersetup{
    colorlinks=true,
    linkcolor=blue!50!black,
    urlcolor=blue!50!black
}

\definecolor{mainblue}{RGB}{33,76,140}
\definecolor{lightblue}{RGB}{235,242,252}
\definecolor{softgray}{RGB}{247,247,247}
\definecolor{linegray}{RGB}{210,210,210}
\definecolor{darkgray}{RGB}{65,65,65}

\setstretch{1.12}
\setlength{\parindent}{0pt}
\setlength{\parskip}{0.45em}

\pagestyle{fancy}
\fancyhf{}
\fancyhead[L]{\textcolor{mainblue}{Operating System Project Documentation}}
\fancyhead[R]{\thepage}
\renewcommand{\headrulewidth}{0.4pt}

\titleformat{\section}
{\Large\bfseries\color{mainblue}}
{\thesection}{1em}{}

\titleformat{\subsection}
{\large\bfseries\color{mainblue}}
{\thesubsection}{1em}{}

\titleformat{\subsubsection}
{\normalsize\bfseries\color{darkgray}}
{}{0em}{}

\newtcolorbox{purposebox}{
    colback=lightblue,
    colframe=mainblue!70,
    arc=2mm,
    boxrule=0.6pt,
    left=8pt,
    right=8pt,
    top=6pt,
    bottom=6pt
}

\newtcolorbox{signaturebox}{
    colback=softgray,
    colframe=linegray,
    arc=1.5mm,
    boxrule=0.4pt,
    left=6pt,
    right=6pt,
    top=5pt,
    bottom=5pt
}

\newcolumntype{Y}{>{\raggedright\arraybackslash}X}

\title{\textbf{\textcolor{mainblue}{Operating System Project Documentation}}}
\author{}
\date{}

\begin{document}

\maketitle
\tableofcontents
\newpage

\section{Overview}

This document describes the classes found in the submitted Operating System project.  
For each class, it includes:
\begin{itemize}[leftmargin=1.5em]
    \item the class purpose
    \item the attributes exactly as found in the code
    \item constructors
    \item method headers with return types
    \item a short English description
\end{itemize}

\section{Class: \texttt{Interpreter}}

\begin{purposebox}
The \texttt{Interpreter} class is responsible for parsing and executing instructions for a process. It uses \texttt{SystemCalls} to interact with input, output, and files, uses \texttt{Memory} to store and retrieve variables, and communicates with the \texttt{OS} when semaphore operations are needed.
\end{purposebox}

\subsection{Attributes}

\rowcolors{2}{softgray}{white}
\begin{tabularx}{\textwidth}{p{4cm} p{3cm} Y}
\toprule
\textbf{Attribute} & \textbf{Type} & \textbf{Description} \\
\midrule
systemCalls & \texttt{SystemCalls} & Stores the system calls object used for input, output, and file operations. \\
memory & \texttt{Memory} & Stores the memory object used to read instructions and variables. \\
operatingSystem & \texttt{OS} & Stores a reference to the operating system so semaphore operations can be delegated to it. \\
\bottomrule
\end{tabularx}

\subsection{Constructors}

\begin{signaturebox}
\texttt{public Interpreter(SystemCalls systemCalls, Memory memory, OS os)}
\end{signaturebox}

Initializes the interpreter with the required \texttt{SystemCalls}, \texttt{Memory}, and \texttt{OS} objects.

\subsection{Methods}

\begin{signaturebox}
\texttt{public void ExecuteInstruction(Process p, String instruction)}
\end{signaturebox}
Parses one instruction and executes it for the given process.

\begin{signaturebox}
\texttt{public void printFromTo(int start, int end)}
\end{signaturebox}
Prints all integer values from \texttt{start} to \texttt{end}, inclusive.

\begin{signaturebox}
\texttt{public void assign()}
\end{signaturebox}
A placeholder method for assignment handling. In the submitted code, it has an empty body.

\begin{signaturebox}
\texttt{public ArrayList<String> readProgramFile(String fileName)}
\end{signaturebox}
Reads a program file line by line and returns its instructions in an \texttt{ArrayList<String>}.

\section{Class: \texttt{Memory}}

\begin{purposebox}
The \texttt{Memory} class simulates RAM. It stores process metadata, instructions, and variables using an array of \texttt{MemoryWord} objects, and provides methods for allocation, access, and variable management.
\end{purposebox}

\subsection{Attributes}

\rowcolors{2}{softgray}{white}
\begin{tabularx}{\textwidth}{p{4cm} p{3cm} Y}
\toprule
\textbf{Attribute} & \textbf{Type} & \textbf{Description} \\
\midrule
SIZE & \texttt{static final int} & Constant memory size set to 40. \\
memory & \texttt{MemoryWord[]} & The array that represents memory cells. \\
\bottomrule
\end{tabularx}

\subsection{Constructors}

\begin{signaturebox}
\texttt{public Memory()}
\end{signaturebox}

Creates a memory object and initializes the array of memory words.

\subsection{Methods}

\begin{signaturebox}
\texttt{public void writeWord(int index, MemoryWord word)}
\end{signaturebox}
Writes a \texttt{MemoryWord} into a specific memory index.

\begin{signaturebox}
\texttt{public MemoryWord readWord(int index)}
\end{signaturebox}
Returns the \texttt{MemoryWord} stored at a specific memory index.

\begin{signaturebox}
\texttt{public void printMemory()}
\end{signaturebox}
Prints the full memory contents to the console.

\begin{signaturebox}
\texttt{public int getRequiredSize(Process process)}
\end{signaturebox}
Calculates how many memory cells are required to store a process.

\begin{signaturebox}
\texttt{public boolean isFreeBlock(int start, int size)}
\end{signaturebox}
Checks whether a block of memory is free starting from a given index.

\begin{signaturebox}
\texttt{public int findFreeBlock(int size)}
\end{signaturebox}
Searches for a contiguous free block of the requested size and returns its start index.

\begin{signaturebox}
\texttt{public boolean allocateProcess(Process process)}
\end{signaturebox}
Allocates memory for a process, stores its PCB information, instructions, and variable slots, and returns whether allocation succeeded.

\begin{signaturebox}
\texttt{public int getVariablesStart(Process process)}
\end{signaturebox}
Returns the starting index of the variable area for the given process.

\begin{signaturebox}
\texttt{public Object getVariableValue(Process process, String variableName)}
\end{signaturebox}
Returns the value of a process variable with the given name.

\begin{signaturebox}
\texttt{public String getInstruction(Process process)}
\end{signaturebox}
Returns the current instruction of a process based on its program counter.

\begin{signaturebox}
\texttt{public void setVariableValue(Process process, String variableName, Object value)}
\end{signaturebox}
Sets or updates a variable value for the given process inside memory.

\section{Class: \texttt{MemoryWord}}

\begin{purposebox}
The \texttt{MemoryWord} class represents one cell in memory. Each memory word stores a name and a value.
\end{purposebox}

\subsection{Attributes}

\rowcolors{2}{softgray}{white}
\begin{tabularx}{\textwidth}{p{4cm} p{3cm} Y}
\toprule
\textbf{Attribute} & \textbf{Type} & \textbf{Description} \\
\midrule
name & \texttt{String} & The label or identifier of the memory word. \\
value & \texttt{Object} & The stored value of the memory word. \\
\bottomrule
\end{tabularx}

\subsection{Constructors}

\begin{signaturebox}
\texttt{public MemoryWord(String name, Object value)}
\end{signaturebox}

Creates a memory word with a name and a value.

\subsection{Methods}

\begin{signaturebox}
\texttt{public String getName()}
\end{signaturebox}
Returns the name of the memory word.

\begin{signaturebox}
\texttt{public void setName(String name)}
\end{signaturebox}
Updates the name of the memory word.

\begin{signaturebox}
\texttt{public Object getValue()}
\end{signaturebox}
Returns the value stored in the memory word.

\begin{signaturebox}
\texttt{public void setValue(Object value)}
\end{signaturebox}
Updates the value stored in the memory word.

\begin{signaturebox}
\texttt{public String toString()}
\end{signaturebox}
Returns a string representation of the memory word.

\section{Class: \texttt{Mutex}}

\begin{purposebox}
The \texttt{Mutex} class simulates a lock for a named shared resource. It tracks whether the resource is locked, which process owns it, and which processes are blocked waiting for it.
\end{purposebox}

\subsection{Attributes}

\rowcolors{2}{softgray}{white}
\begin{tabularx}{\textwidth}{p{4cm} p{3cm} Y}
\toprule
\textbf{Attribute} & \textbf{Type} & \textbf{Description} \\
\midrule
resourceName & \texttt{String} & The name of the protected resource. \\
locked & \texttt{boolean} & Indicates whether the mutex is currently locked. \\
owner & \texttt{Process} & The process that currently owns the mutex. \\
blockedQueue & \texttt{Queue<Process>} & Queue of processes blocked while waiting for the resource. \\
\bottomrule
\end{tabularx}

\subsection{Constructors}

\begin{signaturebox}
\texttt{public Mutex(String resourceName)}
\end{signaturebox}

Creates a mutex for a specific resource and initializes it as unlocked.

\subsection{Methods}

\begin{signaturebox}
\texttt{public boolean semWait(Process process)}
\end{signaturebox}
Tries to acquire the mutex. If it is free, the process becomes the owner; otherwise the process is blocked and added to the blocked queue.

\begin{signaturebox}
\texttt{public Process semSignal(Process process)}
\end{signaturebox}
Releases the mutex if the calling process is the owner. It may unblock and return the next waiting process.

\section{Class: \texttt{OS}}

\begin{purposebox}
The \texttt{OS} class represents the main operating system controller. It initializes the scheduler, memory, interpreter, mutexes, and system calls, and runs the main scheduling and execution loop.
\end{purposebox}

\subsection{Attributes}

\rowcolors{2}{softgray}{white}
\begin{tabularx}{\textwidth}{p{4cm} p{3cm} Y}
\toprule
\textbf{Attribute} & \textbf{Type} & \textbf{Description} \\
\midrule
globalTime & \texttt{int} & Stores the current global time of the simulation. \\
userInput & \texttt{Mutex} & Mutex for the \texttt{userInput} resource. \\
userOutput & \texttt{Mutex} & Mutex for the \texttt{userOutput} resource. \\
file & \texttt{Mutex} & Mutex for the \texttt{file} resource. \\
memory & \texttt{Memory} & The simulated memory used by the operating system. \\
scheduler & \texttt{Scheduler} & The scheduler responsible for selecting processes. \\
sys & \texttt{SystemCalls} & The object used for system call operations. \\
interpreter & \texttt{Interpreter} & The interpreter used to execute instructions. \\
\bottomrule
\end{tabularx}

\subsection{Constructors}

\begin{signaturebox}
\texttt{public OS()}
\end{signaturebox}

Initializes the operating system and creates the memory, scheduler, system calls object, mutexes, and interpreter.

\subsection{Methods}

\begin{signaturebox}
\texttt{public static void main(String[] args)}
\end{signaturebox}
Program entry point. Creates an \texttt{OS} object and starts execution by calling \texttt{run()}.

\begin{signaturebox}
\texttt{public void run()}
\end{signaturebox}
Runs the operating system simulation loop, creates sample processes, loads instructions, allocates memory, schedules processes, and executes instructions.

\begin{signaturebox}
\texttt{public int getGlobalTime()}
\end{signaturebox}
Returns the current value of global time.

\begin{signaturebox}
\texttt{public void incrementGlobalTime()}
\end{signaturebox}
Increments global time by one.

\begin{signaturebox}
\texttt{public void semWait(Process p, String resourceName)}
\end{signaturebox}
Gets the correct mutex by name and calls its \texttt{semWait} operation.

\begin{signaturebox}
\texttt{public void semSignal(Process p, String resourceName)}
\end{signaturebox}
Gets the correct mutex by name and calls its \texttt{semSignal} operation.

\begin{signaturebox}
\texttt{private Mutex getMutexByName(String resourceName)}
\end{signaturebox}
Returns the mutex corresponding to a resource name, or \texttt{null} if the name is invalid.

\section{Class: \texttt{PCB}}

\begin{purposebox}
The \texttt{PCB} class stores the process control block information for a process, including its identifier, state, program counter, and memory boundaries.
\end{purposebox}

\subsection{Attributes}

\rowcolors{2}{softgray}{white}
\begin{tabularx}{\textwidth}{p{4cm} p{3cm} Y}
\toprule
\textbf{Attribute} & \textbf{Type} & \textbf{Description} \\
\midrule
processID & \texttt{int} & The process identifier. \\
processState & \texttt{ProcessState} & The current state of the process. \\
programCounter & \texttt{int} & The index of the next instruction to execute. \\
memStart & \texttt{int} & The starting memory index of the process. \\
memEnd & \texttt{int} & The ending memory index of the process. \\
\bottomrule
\end{tabularx}

\subsection{Constructors}

\begin{signaturebox}
\texttt{public PCB(int pid)}
\end{signaturebox}

Creates a PCB for a process and initializes its state to \texttt{NEW}, its program counter to 0, and its memory boundaries to 0.

\begin{signaturebox}
\texttt{public PCB(int processID, ProcessState processState, int programCounter, int memStart, int memEnd)}
\end{signaturebox}

Creates a PCB using explicit values for all fields.

\subsection{Methods}

\begin{signaturebox}
\texttt{public String toString()}
\end{signaturebox}
Returns a string representation of the PCB.

\section{Class: \texttt{Process}}

\begin{purposebox}
The \texttt{Process} class represents a process in the operating system. It stores its PCB, instructions, arrival time, and status flags such as blocking and whether it is currently in memory.
\end{purposebox}

\subsection{Attributes}

\rowcolors{2}{softgray}{white}
\begin{tabularx}{\textwidth}{p{4cm} p{3cm} Y}
\toprule
\textbf{Attribute} & \textbf{Type} & \textbf{Description} \\
\midrule
pcb & \texttt{PCB} & The PCB associated with the process. \\
instructions & \texttt{ArrayList<String>} & The instruction list of the process. \\
arrivalTime & \texttt{int} & The time at which the process arrives. \\
isBlocked & \texttt{boolean} & Indicates whether the process is blocked. \\
inMemory & \texttt{boolean} & Indicates whether the process is currently loaded in memory. \\
\bottomrule
\end{tabularx}

\subsection{Constructors}

\begin{signaturebox}
\texttt{public Process(int arrivalTime)}
\end{signaturebox}

Creates a process with a given arrival time and initializes its flags.

\subsection{Methods}

\begin{signaturebox}
\texttt{public PCB getPcb()}
\end{signaturebox}
Returns the PCB of the process.

\begin{signaturebox}
\texttt{public void setPcb(PCB pcb)}
\end{signaturebox}
Sets the PCB of the process.

\begin{signaturebox}
\texttt{public ArrayList<String> getInstructions()}
\end{signaturebox}
Returns the instruction list of the process.

\begin{signaturebox}
\texttt{public void setInstructions(ArrayList<String> instructions)}
\end{signaturebox}
Sets the instruction list of the process.

\begin{signaturebox}
\texttt{public int getArrivalTime()}
\end{signaturebox}
Returns the arrival time of the process.

\begin{signaturebox}
\texttt{public void setArrivalTime(int arrivalTime)}
\end{signaturebox}
Updates the arrival time of the process.

\begin{signaturebox}
\texttt{public int getInstructionCounter()}
\end{signaturebox}
Returns the number of instructions in the process.

\begin{signaturebox}
\texttt{public boolean isCompleted()}
\end{signaturebox}
Returns whether the process has finished all its instructions.

\begin{signaturebox}
\texttt{public String toString()}
\end{signaturebox}
Returns a string representation of the process.

\section{Enum: \texttt{ProcessState}}

\begin{purposebox}
The \texttt{ProcessState} enum defines the possible states a process can have during execution.
\end{purposebox}

\subsection{Values}

\begin{itemize}[leftmargin=1.5em]
    \item \texttt{NEW} --- the process has been created but has not started yet
    \item \texttt{READY} --- the process is ready to run
    \item \texttt{RUNNING} --- the process is currently executing
    \item \texttt{BLOCKED} --- the process is waiting for a resource
    \item \texttt{TERMINATED} --- the process has finished execution
\end{itemize}

\section{Class: \texttt{Scheduler}}

\begin{purposebox}
The \texttt{Scheduler} class manages the ready queue and decides which process should run according to the selected scheduling algorithm.
\end{purposebox}

\subsection{Attributes}

\rowcolors{2}{softgray}{white}
\begin{tabularx}{\textwidth}{p{4cm} p{3cm} Y}
\toprule
\textbf{Attribute} & \textbf{Type} & \textbf{Description} \\
\midrule
readyQueue & \texttt{Deque<Process>} & The queue of ready processes. \\
usedTime & \texttt{int} & Tracks the amount of time already used in the current Round Robin quantum. \\
HRRNprocess & \texttt{Process} & Stores the currently selected HRRN process. \\
\bottomrule
\end{tabularx}

\subsection{Constructors}

No explicit constructor is defined, so Java provides the default constructor.

\subsection{Methods}

\begin{signaturebox}
\texttt{public Process roundRobin(int timeQuantum)}
\end{signaturebox}
Chooses the next process using the Round Robin policy.

\begin{signaturebox}
\texttt{public Process HRRN(int globalTime)}
\end{signaturebox}
Chooses the next process using the Highest Response Ratio Next policy.

\begin{signaturebox}
\texttt{public Process MultilevelFeedbackQueue(int globalTime)}
\end{signaturebox}
Placeholder method for a future Multilevel Feedback Queue scheduler.

\begin{signaturebox}
\texttt{public Process SchedulingAlgorithm(String algorithm, int globalTime)}
\end{signaturebox}
Selects which scheduling method to use based on the algorithm name.

\begin{signaturebox}
\texttt{public void addProcess(Process process)}
\end{signaturebox}
Adds a process to the ready queue and sets its state to \texttt{READY}.

\section{Class: \texttt{SystemCalls}}

\begin{purposebox}
The \texttt{SystemCalls} class provides input, output, file operations, and helper methods to read from and write to memory.
\end{purposebox}

\subsection{Attributes}

\rowcolors{2}{softgray}{white}
\begin{tabularx}{\textwidth}{p{4cm} p{3cm} Y}
\toprule
\textbf{Attribute} & \textbf{Type} & \textbf{Description} \\
\midrule
scanner & \texttt{Scanner} & Scanner object used to read user input. \\
\bottomrule
\end{tabularx}

\subsection{Constructors}

\begin{signaturebox}
\texttt{public SystemCalls()}
\end{signaturebox}

Creates a \texttt{SystemCalls} object and initializes the scanner.

\subsection{Methods}

\begin{signaturebox}
\texttt{public void print(String message)}
\end{signaturebox}
Prints a message to the console.

\begin{signaturebox}
\texttt{public String input(String message)}
\end{signaturebox}
Displays a prompt and returns the user's input.

\begin{signaturebox}
\texttt{public void writeFile(String fileName, String data)}
\end{signaturebox}
Writes data into a file.

\begin{signaturebox}
\texttt{public String readFile(String fileName)}
\end{signaturebox}
Reads a file and returns its contents as a string.

\begin{signaturebox}
\texttt{public MemoryWord readFromMemory(Memory memory, int address)}
\end{signaturebox}
Reads and returns a \texttt{MemoryWord} from the given memory address.

\begin{signaturebox}
\texttt{public void writeToMemory(Memory memory, int address, String name, Object value)}
\end{signaturebox}
Creates a new \texttt{MemoryWord} and writes it to the specified memory address.

\section{Class: \texttt{disk}}

\begin{purposebox}
The \texttt{disk} class is a very small placeholder for secondary storage. In the submitted code, it only contains one array attribute and no methods.
\end{purposebox}

\subsection{Attributes}

\rowcolors{2}{softgray}{white}
\begin{tabularx}{\textwidth}{p{4cm} p{3cm} Y}
\toprule
\textbf{Attribute} & \textbf{Type} & \textbf{Description} \\
\midrule
disk & \texttt{Object[]} & Array representing the simulated disk contents. Its size is 100. \\
\bottomrule
\end{tabularx}

\subsection{Constructors}

No explicit constructor is defined, so Java provides the default constructor.

\subsection{Methods}

No methods are defined in this class.

\section{Notes}

This documentation matches the uploaded code much more closely than the previous version.  
It documents what is actually present in the source, even when a method is unfinished, placeholder-only, or partially implemented.

\end{document}